\section*{Bab \ref{ch:g6:fractions} --- Pecahan: Langkah Pertama}

\begin{solution}{exo:g6:fractions:1}
Kue: $\frac38$ dimakan. \quad Cokelat: $\frac{7}{12}$. \quad
Pai: $\frac66 = 1$ (seluruh painya).
\end{solution}

\begin{solution}{exo:g6:fractions:2}
Empat bagian dari lima diberi warna. $\frac45 < 1$ (empat bagian dari
lima kurang daripada seluruh batangnya); $\frac65 > 1$ (enam perlima
adalah satu batang utuh ditambah seperlima lagi).
\end{solution}

\begin{solution}{exo:g6:fractions:3}
Pada garis berskala sepertigaan: $\frac13$ berada satu tanda sesudah
$0$; $\frac53$ berada satu tanda sebelum $2$ (karena $2 = \frac63$);
$\frac73$ berada satu tanda sesudah $2$. Jadi $\frac53$ dan $\frac73$
sama-sama berjarak $\frac13$ dari $2$.
\end{solution}

\begin{solution}{exo:g6:fractions:4}
$\frac94 = 2 + \frac14$; \quad
$\frac{13}{5} = 2 + \frac35$; \quad
$\frac{12}{6} = 2$ (tepat); \quad
$\frac{25}{8} = 3 + \frac18$.
\end{solution}

\begin{solution}{exo:g6:fractions:5}
$\frac23 = \frac{8}{12}$ (kalikan dengan $4$); \quad
$\frac{15}{20} = \frac34$ (bagilah dengan $5$); \quad
$\frac{7}{10} = \frac{21}{30}$ (kalikan dengan $3$); \quad
$\frac56 = \frac{10}{12}$ (bagilah dengan $2$).
\end{solution}

\begin{solution}{exo:g6:fractions:6}
$\frac68 = \frac34$; \quad
$\frac{15}{25} = \frac35$; \quad
$\frac{18}{24} = \frac{9}{12} = \frac34$; \quad
$\frac{35}{7} = 5$ (sebuah bilangan bulat).
\end{solution}

\begin{solution}{exo:g6:fractions:7}
$\frac12$ dari $86$: $86 \div 2 = 43$.

$\frac34$ dari $60$: $60 \div 4 = 15$, lalu $15 \times 3 = 45$.

$\frac25$ dari $35$: $35 \div 5 = 7$, lalu $7 \times 2 = 14$.

$\frac{7}{10}$ dari $250$: $250 \div 10 = 25$, lalu $25 \times 7 = 175$.
\end{solution}

\begin{solution}{exo:g6:fractions:8}
$\frac12 > \frac13$: setengahan adalah bagian yang lebih besar daripada sepertigaan.
$\frac35 > \frac25$: bagian yang sama, tetapi lebih banyak.
$\frac34 = \frac68 < \frac78$: tujuh perdelapan mengalahkan enam perdelapan.
\end{solution}

\begin{solution}{exo:g6:fractions:9}
$\frac12 = 0.5$; \quad $\frac34 = 0.75$; \quad
$\frac{7}{10} = 0.7$; \quad $\frac94 = 2.25$. Yang berbeda sendiri
adalah $\frac13 = 0.333\dots$, yang penulisan desimalnya tidak pernah
berakhir: ia bukan \omterm{def:g6:decimals:places}{bilangan desimal}.
\end{solution}

\begin{solution}{exo:g6:fractions:10}
Yang naik bus: $\frac34$ dari $28$: $28 \div 4 = 7$, lalu
$7 \times 3 = 21$ murid. Yang berjalan kaki: $28 - 21 = 7$ murid
(itulah \omterm{def:g3:division:half}{seperempat} sisanya: $28 \div 4 = 7$). Periksa: $21 + 7 = 28$.
\end{solution}

\begin{solution}{exo:g6:fractions:11}
Dua pertiga tabungannya adalah $18$ euro, jadi \emph{satu} pertiganya
$18 \div 2 = 9$ euro, dan seluruh tabungannya (tiga pertiga) adalah
$9 \times 3 = 27$ euro. Periksa: $\frac23$ dari $27$ adalah $18$.
\end{solution}

\begin{solution}{exo:g6:fractions:12}
Air yang terpakai turun dari $\frac58$ menjadi $\frac48$ (setengah
sama dengan empat perdelapan) isi tangki: $10$ liter itu mewakili
$\frac58 - \frac48 = \frac18$ kapasitasnya. Jadi tangki yang penuh
memuat $10 \times 8 = 80$ liter. Periksa: $\frac58$ dari $80$ adalah
$50$, dikurangi $10$ menyisakan $40$, yang memang setengah dari $80$.
\end{solution}

\begin{solution}{pb:g6:fractions:1}
\textbf{1.} Setengah dari $16$ kotak adalah $8$ kotak. Sebagai
\omterm{def:g6:fractions:def}{pecahan}: $\frac12$ batangnya, dan karena setiap setengahnya $8$ kotak
dari $16$, juga $\frac{8}{16}$ --- \omterm{def:g6:fractions:def}{pecahan} yang sama ditulis dengan
dua cara (\cref{prop:g6:fractions:equal}).

\textbf{2.} Setengah dari $8$ kotak yang tersisa adalah $4$ kotak,
yaitu $\frac{4}{16}$ batangnya. Disederhanakan dengan $4$:
$\frac{4}{16} = \frac14$. Pada gambarnya: memotong setengah yang
tersisa menjadi dua akan menghasilkan dua \omterm{def:g3:division:half}{seperempat} dari batang
semula --- setengah dari setengah adalah \omterm{def:g3:division:half}{seperempat}.

\textbf{3.} Gigitan ketiga: setengah dari $4$ kotak adalah $2$ kotak,
yaitu $\frac{2}{16} = \frac18$ batangnya. Gigitan keempat: $1$ kotak,
yaitu $\frac{1}{16}$. Dari satu gigitan ke gigitan berikutnya,
\omterm{def:g4:fractions:def}{penyebutnya} \emph{berlipat dua}: $2$, $4$, $8$, $16$.

\textbf{4.} Yang dimakan: $8 + 4 + 2 + 1 = 15$ kotak dari $16$, yaitu
$\frac{15}{16}$ batangnya. Yang tersisa: $1$ kotak, $\frac{1}{16}$.

\textbf{5.} $\frac12 = \frac{8}{16}$ (kalikan atas dan bawahnya dengan
$8$), $\frac14 = \frac{4}{16}$ (dengan $4$), $\frac18 = \frac{2}{16}$
(dengan $2$), dan $\frac{1}{16}$ tetap. Membilang perenambelasannya:
ada $8 + 4 + 2 + 1 = 15$, jadi \omterm{def:g1:addition:def}{jumlahnya} $\frac{15}{16}$.

\textbf{6.} Gigitannya: $32$, $16$, $8$, $4$, $2$, $1$ kotak. Yang
dimakan seluruhnya: $32 + 16 + 8 + 4 + 2 + 1 = 63$ kotak, jadi $1$
kotak tersisa: $\frac{1}{64}$ batangnya.

\textbf{7.} Setelah setiap gigitan yang tersisa adalah
\[
\frac12,\quad \frac14,\quad \frac18,\quad \frac{1}{16},\quad
\frac{1}{32},\quad \frac{1}{64} :
\]
sisanya membelah dua setiap kali --- \omterm{def:g4:fractions:def}{penyebutnya} berlipat dua.
Gigitan ketujuh yang khayalan itu akan menyisakan $\frac{1}{128}$.

\textbf{8.} Setiap gigitan hanya memakan \emph{setengah} dari yang
tersisa, jadi setengah sisanya masih ada di sana: setelah berapa pun
gigitan, selalu ada yang tertinggal. Karena itu \omterm{def:g1:addition:def}{jumlah} yang dimakan
selalu di bawah $1$ --- persis seperti dua puluh angka sembilan pada
\cref{pb:g6:decimals:1}, yang makin lama makin dekat ke $3$ tanpa
pernah mencapainya.

\textbf{9.} Kedua \omterm{def:g6:fractions:def}{pecahan} itu berpembilang $1$, dan memotong menjadi
$128$ bagian membuat bagian yang lebih kecil daripada memotong menjadi
$100$: jadi $\frac{1}{128} < \frac{1}{100}$. Setelah gigitan ketujuh
\omterm{def:g3:division:remainder}{sisanya} $\frac{1}{128}$, jadi bagian yang dimakan melebihi
$1 - \frac{1}{100} = \frac{99}{100}$: tantangan itu dimenangkan pada
gigitan ketujuh.

\textbf{10.} Dalam perenambelasan, \omterm{def:g1:addition:def}{jumlahnya} adalah $\frac{8}{16}$,
$\frac{12}{16}$, $\frac{14}{16}$, $\frac{15}{16}$. Setiap \omterm{def:g1:addition:def}{jumlah} baru
mendarat tepat \emph{di tengah} antara \omterm{def:g1:addition:def}{jumlah} sebelumnya dan $1$:
tangganya terus membelah dua \omterm{def:g3:division:remainder}{sisa} jaraknya ke $1$, tanpa pernah
menginjaknya.

\textbf{11.} Yang benar: pada setiap tahap gambaran Zeno, memang masih
ada jarak yang tersisa (pertanyaan 8) --- daftar tahapnya tidak pernah
berakhir. Perangkapnya: tahapnya menjadi sangat pendek, dalam jarak
\emph{dan} dalam waktu berlari; pelari itu tidak menghabiskan waktu
yang sama pada setiap tahap. Menjelaskan lariannya dalam potongan yang
tak hingga banyaknya dan makin mengecil tidak membuat larian itu
sendiri tanpa akhir: pelarinya sampai ke tembok, dan pertanyaan 9
menunjukkan gambarannya melewati titik mana pun sebelum tembok itu.

\textbf{12.} Potonglah $8$ batang menjadi setengahan: $16$ potongan
setengah. Potonglah $4$ batang menjadi seperempatan: $16$ potongan
\omterm{def:g3:division:half}{seperempat}. Potonglah $2$ batang menjadi perdelapanan: $16$ potongan
perdelapan. Potonglah batang terakhir menjadi perenambelasan: $16$
potongan kecil. Itu memakai $8 + 4 + 2 + 1 = 15$ batang, dan setiap
anak menerima satu potongan dari tiap ukuran.

\textbf{13.} Bagian setiap anak adalah
\[
\frac12 + \frac14 + \frac18 + \frac{1}{16} = \frac{15}{16}
\]
sebuah batang (pertanyaan 5). Enam belas bagian semacam itu menjadi
$16 \times \frac{15}{16} = 15$ batang
(\cref{thm:g6:fractions:quotient}): tepat seperti yang dipotong, tanpa
\omterm{def:g3:division:remainder}{sisa} sama sekali --- pembagian $15$ batang yang adil kepada $16$ anak,
hanya dengan membelah dua.

\textbf{14.} Tetangga sebelah menerima $\frac78$ batang untuk
masing-masing. Dengan \omterm{def:g4:fractions:def}{penyebut} $16$: $\frac78 = \frac{14}{16}$,
sedangkan anak kita menerima $\frac{15}{16}$. Jadi satu anak dari yang
$16$ menerima lebih banyak: $\frac{15}{16} > \frac{14}{16}$, \omterm{ex:g1:subtraction:difference}{selisih}
seperenam belas batang.

\textbf{15.} Pada gambar persegi itu, setiap gigitan yang diwarnai
mengisi setengah dari pojok yang belum terwarnai, dan pojok yang belum
terwarnai itu terus mengecil: setelah setiap gigitan ia setengah dari
sebelumnya, dan ia tidak pernah lenyap. Dua fakta yang benar: (i)
setelah gigitan yang banyaknya berhingga, selalu ada \omterm{def:g3:division:remainder}{sisa} --- \omterm{def:g1:addition:def}{jumlah}
yang dimakan selalu tepat di bawah $1$; (ii) \omterm{def:g3:division:remainder}{sisa} itu menjadi lebih
kecil daripada \omterm{def:g6:fractions:def}{pecahan} mana pun yang ingin disebut orang
($\frac{1}{100}$, $\frac{1}{1000}$, \dots), asalkan ia menggigit cukup
lama. Memberi makna yang tepat pada kata ``\omterm{def:g1:addition:def}{jumlah} tak hingga itu sama
dengan $1$'' adalah tugas \emph{limit}, di buku Sekolah Menengah.

\end{solution}
